% 前言
\frontsection{前言}

这份笔记记录作者在山东大学 2026 年秋季学期修读“高等概率论”课程的内容。课程以 Durrett 的 \emph{Probability: Theory and Examples}（第五版）为教材，内容大致覆盖该书的前五章：测度论基础、大数定律、中心极限定理、鞅与 Markov 链。

笔记并不逐次记录课堂，而是按照理论本身的逻辑顺序重新组织。课堂上讲过的内容都在这里；有些地方课堂上一带而过，这里会补全必要的背景与完整的证明。写作时尽量做到自成体系：除了数学分析与线性代数的基本知识外，不假定读者学过实变函数或测度论。凡是后文要用到的测度论结果，都在正文里从定义出发给出证明。

\subsection*{笔记的体例}

体例仿照于品《数学分析之课程讲义》。每一讲是一个相对完整的主题，标题罗列本讲建立的知识点。正文的重点是把一套理论讲清楚：为什么要引入某个概念，定理的每个条件起什么作用，证明是怎样一步一步走通的。定义、命题、引理、定理与推论共用一个全书连续的编号；例子、注记与练习不编号。例子是正文的一部分，用来说明概念的含义、定理的用法，以及去掉某个条件之后结论如何失效。

证明一律写完整。测度论与概率论里的许多论证，难点不在想法，而在细节：可测性是否成立，极限与积分能否交换，某个集合是否真的属于给定的 $\sigma$-代数。这些地方笔记都不省略。

\subsection*{作业}

每讲之后附有作业，分为三部分：
\begin{enumerate}
  \item \textbf{基本习题}：课堂内容的补充，以及概念与方法的训练；
  \item \textbf{思考题}：难度较大或者延伸较远的问题，难题标有 *；
  \item \textbf{课外知识}：课堂之外的一个完整结果。
\end{enumerate}
较大的结果采用“法式大题”的形式给出：把一个定理的证明拆成若干相互衔接的小问，前一问构造的对象、证明的引理或者完成的计算，是后一问的台阶。沿着这样的线索做下去，既能熟悉基本的论证技术，也能看清一个完整的证明是怎样搭起来的。每道题后面都附有详细的解答，可以先自己想，想不出来时再读解答，学习其中的方法。
